\iftestbuild % TEST-ONLY-BEGIN
\setcounter{section}{2}\setcounter{theorem}{33}
\fi % TEST-ONLY-END
% Source: book pp. 44--50 (PDF 58--64); solutions pp. 42--53, 169.
% A numbered display (Axler's 2.44, 2.45) takes its number from the shared
% item counter: \axeqnum{2.44} inside equation* (also defined in 2B.tex).
\providecommand\axeqnum[1]{\refstepcounter{theorem}\tag*{\thetheorem}\label{ax:#1}}
\section{Dimension}

Although we have been discussing finite-dimensional vector spaces, we have not
yet defined the dimension of such an object. How should dimension be defined? A
reasonable definition should force the dimension of $\F^n$ to equal $n$. Notice
that the standard basis
\[ (1,0,\dots,0),\ (0,1,0,\dots,0),\ \dots,\ (0,\dots,0,1) \]
of $\F^n$ has length $n$. Thus we are tempted to define the dimension as the
length of a basis. However, a finite-dimensional vector space in general has
many different bases, and our attempted definition makes sense only if all bases
in a given vector space have the same length. Fortunately that turns out to be
the case, as we now show.

\begin{result}[Basis length does not depend on basis]\label{ax:2.34}
Any two bases of a finite-dimensional vector space have the same length.
\end{result}

\begin{proof}
Suppose $V$ is finite-dimensional. Let $B_1$ and $B_2$ be two bases of $V$. Then
$B_1$ is linearly independent in $V$ and $B_2$ spans $V$, so the length of $B_1$
is at most the length of $B_2$ (by \ref{ax:2.22}). Interchanging the roles of
$B_1$ and $B_2$, we also see that the length of $B_2$ is at most the length of
$B_1$. Thus the length of $B_1$ equals the length of $B_2$, as desired.
\end{proof}

Now that we know that any two bases of a finite-dimensional vector space have
the same length, we can formally define the dimension of such spaces.

\begin{definition}[Dimension, $\dim V$]\label{ax:2.35}\leavevmode
\begin{itemize}
\item The \emph{dimension} of a finite-dimensional vector space is the length
  of any basis of the vector space.
\item The dimension of a finite-dimensional vector space $V$ is denoted by
  $\dim V$.
\end{itemize}
\end{definition}

\begin{example}[Dimensions]\label{ax:2.36}\leavevmode
\begin{itemize}
\item $\dim\F^n=n$ because the standard basis of $\F^n$ has length $n$.
\item $\dim\Poly_m(\F)=m+1$ because the standard basis $1,z,\dots,z^m$ of
  $\Poly_m(\F)$ has length $m+1$.
\item If $U=\{(x,x,y)\in\F^3 : x,y\in\F\}$, then $\dim U=2$ because
  $(1,1,0),(0,0,1)$ is a basis of $U$.
\item If $U=\{(x,y,z)\in\F^3 : x+y+z=0\}$, then $\dim U=2$ because the list
  $(1,-1,0),(1,0,-1)$ is a basis of $U$.
\end{itemize}
\end{example}

Every subspace of a finite-dimensional vector space is finite-dimensional
(by \ref{ax:2.25}) and so has a dimension. The next result gives the expected
inequality about the dimension of a subspace.

\begin{result}[Dimension of a subspace]\label{ax:2.37}
If $V$ is finite-dimensional and $U$ is a subspace of $V$, then
$\dim U\le\dim V$.
\end{result}

\begin{proof}
Suppose $V$ is finite-dimensional and $U$ is a subspace of $V$. Think of a basis
of $U$ as a linearly independent list in $V$, and think of a basis of $V$ as a
spanning list in $V$. Now use \ref{ax:2.22} to conclude that
$\dim U\le\dim V$.
\end{proof}

\marginnote[-4mm]{The real vector space $\R^2$ has dimension two; the complex vector
space $\C$ has dimension one. As sets, $\R^2$ can be identified with $\C$ (and
addition is the same on both spaces, as is scalar multiplication by real
numbers). Thus when we talk about the dimension of a vector space, the role
played by the choice of $\F$ cannot be neglected.}%
To check that a list of vectors in $V$ is a basis of $V$, we must, according to
the definition, show that the list in question satisfies two properties: it
must be linearly independent and it must span $V$. The next two results show
that if the list in question has the right length, then we only need to check
that it satisfies one of the two required properties. First we prove that every
linearly independent list of the right length is a basis.

\begin{result}[Linearly independent list of the right length is a basis]%
\label{ax:2.38}
Suppose $V$ is finite-dimensional. Then every linearly independent list of
vectors in $V$ of length $\dim V$ is a basis of $V$.
\end{result}

\begin{proof}
Suppose $\dim V=n$ and $v_1,\dots,v_n$ is linearly independent in $V$. The list
$v_1,\dots,v_n$ can be extended to a basis of $V$ (by \ref{ax:2.32}). However,
every basis of $V$ has length $n$, so in this case the extension is the trivial
one, meaning that no elements are adjoined to $v_1,\dots,v_n$. Thus
$v_1,\dots,v_n$ is a basis of $V$, as desired.
\end{proof}

The next result is a useful consequence of the previous result.

\begin{result}[Subspace of full dimension equals the whole space]\label{ax:2.39}
Suppose that $V$ is finite-dimensional and $U$ is a subspace of $V$ such that
$\dim U=\dim V$. Then $U=V$.
\end{result}

\begin{proof}
Let $u_1,\dots,u_n$ be a basis of $U$. Thus $n=\dim U$, and by hypothesis we
also have $n=\dim V$. Thus $u_1,\dots,u_n$ is a linearly independent list of
vectors in $V$ (because it is a basis of $U$) of length $\dim V$. From
\ref{ax:2.38}, we see that $u_1,\dots,u_n$ is a basis of $V$. In particular
every vector in $V$ is a linear combination of $u_1,\dots,u_n$. Thus $U=V$.
\end{proof}

\begin{example}[A basis of $\F^2$]\label{ax:2.40}
Consider the list $(5,7),(4,3)$ of vectors in $\F^2$. This list of length two
is linearly independent in $\F^2$ (because neither vector is a scalar multiple
of the other). Note that $\F^2$ has dimension two. Thus \ref{ax:2.38} implies
that the linearly independent list $(5,7),(4,3)$ of length two is a basis of
$\F^2$ (we do not need to bother checking that it spans $\F^2$).
\end{example}

\begin{example}[A basis of a subspace of $\Poly_3(\R)$]\label{ax:2.41}
Let $U$ be the subspace of $\Poly_3(\R)$ defined by
\[ U = \{p\in\Poly_3(\R) : p'(5)=0\}. \]
To find a basis of $U$, first note that each of the polynomials $1$, $(x-5)^2$,
and $(x-5)^3$ is in $U$.

Suppose $a,b,c\in\R$ and
\[ a+b(x-5)^2+c(x-5)^3 = 0 \]
for every $x\in\R$. Without explicitly expanding the left side of the equation
above, we can see that the left side has a $cx^3$ term. Because the right side
has no $x^3$ term, this implies that $c=0$. Because $c=0$, we see that the left
side has a $bx^2$ term, which implies that $b=0$. Because $b=c=0$, we can also
conclude that $a=0$. Thus the equation above implies that $a=b=c=0$. Hence the
list $1,(x-5)^2,(x-5)^3$ is linearly independent in $U$. Thus $3\le\dim U$.
Hence
\[ 3 \le \dim U \le \dim\Poly_3(\R) = 4, \]
where we have used \ref{ax:2.37}.

The polynomial $x$ is not in $U$ because its derivative is the constant
function $1$. Thus $U\ne\Poly_3(\R)$. Hence $\dim U\ne4$ (by \ref{ax:2.39}). The
inequality above now implies that $\dim U=3$. Thus the linearly independent list
$1,(x-5)^2,(x-5)^3$ in $U$ has length $\dim U$ and hence is a basis of $U$
(by \ref{ax:2.38}).
\end{example}

Now we prove that a spanning list of the right length is a basis.

\begin{result}[Spanning list of the right length is a basis]\label{ax:2.42}
Suppose $V$ is finite-dimensional. Then every list of vectors in $V$ that spans
$V$ and has length $\dim V$ is a basis of $V$.
\end{result}

\begin{proof}
Suppose $\dim V=n$ and $v_1,\dots,v_n$ spans $V$. The list $v_1,\dots,v_n$ can
be reduced to a basis of $V$ (by \ref{ax:2.30}). However, every basis of $V$ has
length $n$, so in this case the reduction is the trivial one, meaning that no
elements are deleted from $v_1,\dots,v_n$. Thus $v_1,\dots,v_n$ is a basis of
$V$, as desired.
\end{proof}

The next result gives a formula for the dimension of the sum of two subspaces
of a finite-dimensional vector space. This formula is analogous to a familiar
counting formula: the number of elements in the union of two finite sets equals
the number of elements in the first set, plus the number of elements in the
second set, minus the number of elements in the intersection of the two sets.

\begin{result}[Dimension of a sum]\label{ax:2.43}
If $V_1$ and $V_2$ are subspaces of a finite-dimensional vector space, then
\[ \dim(V_1+V_2) = \dim V_1+\dim V_2-\dim(V_1\cap V_2). \]
\end{result}

\begin{proof}
Let $v_1,\dots,v_m$ be a basis of $V_1\cap V_2$; thus $\dim(V_1\cap V_2)=m$.
Because $v_1,\dots,v_m$ is a basis of $V_1\cap V_2$, it is linearly independent
in $V_1$. Hence this list can be extended to a basis
$v_1,\dots,v_m,u_1,\dots,u_j$ of $V_1$ (by \ref{ax:2.32}). Thus
$\dim V_1=m+j$. Also extend $v_1,\dots,v_m$ to a basis
$v_1,\dots,v_m,w_1,\dots,w_k$ of $V_2$; thus $\dim V_2=m+k$.

We will show that
\begin{equation*}
v_1,\dots,v_m,u_1,\dots,u_j,w_1,\dots,w_k \axeqnum{2.44}
\end{equation*}
is a basis of $V_1+V_2$. This will complete the proof, because then we will have
\begin{align*}
\dim(V_1+V_2) &= m+j+k\\
  &= (m+j)+(m+k)-m\\
  &= \dim V_1+\dim V_2-\dim(V_1\cap V_2).
\end{align*}

The list \ref{ax:2.44} is contained in $V_1\cup V_2$ and thus is contained in
$V_1+V_2$. The span of this list contains $V_1$ and contains $V_2$ and hence is
equal to $V_1+V_2$. Thus to show that \ref{ax:2.44} is a basis of $V_1+V_2$ we
only need to show that it is linearly independent.

To prove that \ref{ax:2.44} is linearly independent, suppose
\[ a_1v_1+\cdots+a_mv_m+b_1u_1+\cdots+b_ju_j+c_1w_1+\cdots+c_kw_k = 0, \]
where all the $a$'s, $b$'s, and $c$'s are scalars. We need to prove that all the
$a$'s, $b$'s, and $c$'s equal $0$. The equation above can be rewritten as
\begin{equation*}
c_1w_1+\cdots+c_kw_k = -a_1v_1-\cdots-a_mv_m-b_1u_1-\cdots-b_ju_j,
\axeqnum{2.45}
\end{equation*}
which shows that $c_1w_1+\cdots+c_kw_k\in V_1$. All the $w$'s are in $V_2$, so
this implies that $c_1w_1+\cdots+c_kw_k\in V_1\cap V_2$. Because
$v_1,\dots,v_m$ is a basis of $V_1\cap V_2$, we have
\[ c_1w_1+\cdots+c_kw_k = d_1v_1+\cdots+d_mv_m \]
for some scalars $d_1,\dots,d_m$. But $v_1,\dots,v_m,w_1,\dots,w_k$ is linearly
independent, so the last equation implies that all the $c$'s (and $d$'s)
equal $0$. Thus \ref{ax:2.45} becomes the equation
\[ a_1v_1+\cdots+a_mv_m+b_1u_1+\cdots+b_ju_j = 0. \]
Because the list $v_1,\dots,v_m,u_1,\dots,u_j$ is linearly independent, this
equation implies that all the $a$'s and $b$'s are $0$, completing the proof.
\end{proof}

For $S$ a finite set, let $\#S$ denote the number of elements of $S$. The table
below compares finite sets with finite-dimensional vector spaces, showing the
analogy between $\#S$ (for sets) and $\dim V$ (for vector spaces), as well as
the analogy between unions of subsets (in the context of sets) and sums of
subspaces (in the context of vector spaces).

\begin{center}\small
\begin{tabular}{|p{.44\linewidth}|p{.44\linewidth}|}
\hline
\multicolumn{1}{|c|}{\textbf{sets}} &
\multicolumn{1}{c|}{\textbf{vector spaces}}\\
\hline
$S$ is a finite set & $V$ is a finite-dimensional vector space\\
\hline
$\#S$ & $\dim V$\\
\hline
for subsets $S_1,S_2$ of $S$, the union $S_1\cup S_2$ is the smallest subset of
$S$ containing $S_1$ and $S_2$ &
for subspaces $V_1,V_2$ of $V$, the sum $V_1+V_2$ is the smallest subspace of
$V$ containing $V_1$ and $V_2$\\
\hline
$\#(S_1\cup S_2)$\newline $=\#S_1+\#S_2-\#(S_1\cap S_2)$ &
$\dim(V_1+V_2)$\newline $=\dim V_1+\dim V_2-\dim(V_1\cap V_2)$\\
\hline
$\#(S_1\cup S_2)=\#S_1+\#S_2$\newline $\iff S_1\cap S_2=\varnothing$ &
$\dim(V_1+V_2)=\dim V_1+\dim V_2$\newline $\iff V_1\cap V_2=\{0\}$\\
\hline
$S_1\cup\cdots\cup S_m$ is a disjoint union $\iff$\newline
$\#(S_1\cup\cdots\cup S_m)=\#S_1+\cdots+\#S_m$ &
$V_1+\cdots+V_m$ is a direct sum $\iff$\newline
$\dim(V_1+\cdots+V_m)$\newline $=\dim V_1+\cdots+\dim V_m$\\
\hline
\end{tabular}
\end{center}

The last row above focuses on the analogy between disjoint unions (for sets) and
direct sums (for vector spaces). The proof of the result in the last box above
will be given in \ref{ax:3.94}.

You should be able to find results about sets that correspond, via analogy, to
the results about vector spaces in Exercises 12 through 18.

% ------------------------------------------------------------------ exercises
\begin{exc}

\begin{exercise}
Show that the subspaces of $\R^2$ are precisely $\{0\}$, all lines in $\R^2$
containing the origin, and $\R^2$.
\end{exercise}
\begin{solution}
\solnote{Manual Remark}{In Exercise~2C3 through Exercise~2C7, each problem
gives a set $U$ that satisfies some linear condition on the polynomials in
$\Poly_4(\F)$. In each case, the polynomials we exhibit span $U$, since
expressing a general element of $U$ in terms of the free parameters allowed by
the defining condition writes it directly as a linear combination of them. They
are also linearly independent: since they have different degrees, a linear
combination of them can equal the zero polynomial only if every coefficient is
zero, as otherwise the highest-degree term present could not cancel. We will
therefore omit this reasoning going forward and simply exhibit a basis for each
such $U$.}

Since $\dim(\R^2)=2$, any subspace of $\R^2$ must have dimension $0$, $1$, or
$2$. The only subspace of dimension $0$ is $\{0\}$, and the only subspace of
dimension $2$ is $\R^2$. Any subspace with dimension $1$ consists of all scalar
multiples of some nonzero $(x,y)\in\R^2$, which is precisely the line with
direction $(x,y)$ containing the origin. Thus, the subspaces of $\R^2$ are
precisely $\{0\}$, all lines in $\R^2$ containing the origin, and $\R^2$.
\end{solution}

\begin{exercise}
Show that the subspaces of $\R^3$ are precisely $\{0\}$, all lines in $\R^3$
containing the origin, all planes in $\R^3$ containing the origin, and $\R^3$.
\end{exercise}
\begin{solution}
Since $\dim(\R^3)=3$, any subspace of $\R^3$ must have dimension $0$, $1$, $2$,
or $3$. The only subspace of dimension $0$ is $\{0\}$, and the only subspace of
dimension $3$ is $\R^3$. Any subspace with dimension $1$ consists of all scalar
multiples of some nonzero $(x,y,z)\in\R^3$, which is precisely the line
containing the origin in the direction of $(x,y,z)$. Any subspace with
dimension $2$ consists of all linear combinations of two linearly independent
vectors in $\R^3$, which is precisely a plane containing the origin. Thus, the
subspaces of $\R^3$ are precisely $\{0\}$, all lines in $\R^3$ containing the
origin, all planes in $\R^3$ containing the origin, and $\R^3$.
\end{solution}

\begin{exercise}\leavevmode
\begin{parts}
\item Let $U=\{p\in\Poly_4(\F) : p(6)=0\}$. Find a basis of $U$.
\item Extend the basis in (a) to a basis of $\Poly_4(\F)$.
\item Find a subspace $W$ of $\Poly_4(\F)$ such that $\Poly_4(\F)=U\oplus W$.
\end{parts}
\end{exercise}
\begin{solution}
\solnote{Manual Remark}{Please see remark at the beginning of this section
regarding the basis of $U$.}
\begin{parts}
\item Let $p=a_0+a_1x+a_2x^2+a_3x^3+a_4x^4\in U$. Then we have $p(6)=0$, which
  gives us the equation
  \[ a_0+6a_1+36a_2+216a_3+1296a_4 = 0. \]
  We can solve this equation for $a_0$ in terms of $a_1,a_2,a_3,a_4$:
  \[ a_0 = -6a_1-36a_2-216a_3-1296a_4. \]
  Then we can write any polynomial in $U$ as a linear combination of the
  following four polynomials:
  \begin{gather*}
  u_1=-6+x, \quad u_2=-36+x^2,\\
  u_3=-216+x^3, \quad u_4=-1296+x^4.
  \end{gather*}
\item We can extend the basis $\{u_1,u_2,u_3,u_4\}$ of $U$ to a basis of
  $\Poly_4(\F)$ by adding one more polynomial that is linearly independent of
  $u_1,u_2,u_3,u_4$. For example, we can add the polynomial
  \[ u_5 = 1, \]
  which is linearly independent of $u_1,u_2,u_3,u_4$: since $u_5$ has no
  higher-degree terms, it cannot be expressed as a linear combination of
  $u_1,u_2,u_3,u_4$. Thus, $\{u_1,u_2,u_3,u_4,u_5\}$ is a basis of
  $\Poly_4(\F)$.
\item Let $W=\spn u_5=\spn 1$. Then we have
  \[ \Poly_4(\F) = U+W, \]
  since every polynomial in $\Poly_4(\F)$ can be written as a sum of a
  polynomial in $U$ and a polynomial in $W$. Moreover, if $p\in U\cap W$, then
  $p=\alpha$ for some $\alpha\in\F$. Since $p(6)=0$, we have $\alpha=0$. Thus,
  $U\cap W=\{0\}$, and so we have
  \[ \Poly_4(\F) = U\oplus W. \qedhere \]
\end{parts}
\end{solution}

\begin{exercise}\leavevmode
\begin{parts}
\item Let $U=\{p\in\Poly_4(\R) : p''(6)=0\}$. Find a basis of $U$.
\item Extend the basis in (a) to a basis of $\Poly_4(\R)$.
\item Find a subspace $W$ of $\Poly_4(\R)$ such that $\Poly_4(\R)=U\oplus W$.
\end{parts}
\end{exercise}
\begin{solution}
\solnote{Manual Remark}{Please see remark at the beginning of this section
regarding the basis of $U$.}
\begin{parts}
\item Let $p=a_0+a_1x+a_2x^2+a_3x^3+a_4x^4\in U$. Then we have
  \[ p''(x) = 2a_2+6a_3x+12a_4x^2, \]
  and so $p''(6)=0$ gives us the equation
  \[ 2a_2+36a_3+432a_4 = 0. \]
  We can solve this equation for $a_2$ in terms of $a_3,a_4$:
  \[ a_2 = -18a_3-216a_4. \]
  Then we can write any polynomial in $U$ as a linear combination of the
  following four polynomials:
  \begin{gather*}
  u_1=1, \quad u_2=x,\\
  u_3=-18x^2+x^3, \quad u_4=-216x^2+x^4.
  \end{gather*}
\item We can extend the basis $\{u_1,u_2,u_3,u_4\}$ of $U$ to a basis of
  $\Poly_4(\R)$ by adding one more polynomial that is linearly independent of
  $u_1,u_2,u_3,u_4$. For example, we can add the polynomial
  \[ u_5 = x^2, \]
  which is linearly independent of $u_1,u_2,u_3,u_4$: since only $u_3$ and
  $u_4$ have $x^2$ terms but contain other higher-degree terms, $u_5$ cannot be
  expressed as a linear combination of $u_1,u_2,u_3,u_4$. Thus,
  $\{u_1,u_2,u_3,u_4,u_5\}$ is a basis of $\Poly_4(\R)$.
\item Let $W=\spn u_5=\spn x^2$. Then we have
  \[ \Poly_4(\R) = U+W, \]
  since every polynomial in $\Poly_4(\R)$ can be written as a sum of a
  polynomial in $U$ and a polynomial in $W$. Moreover, if $p\in U\cap W$, then
  $p=\alpha x^2$ for some $\alpha\in\R$. Since $p''(6)=0$, we have $2\alpha=0$,
  which implies that $\alpha=0$. Thus, $U\cap W=\{0\}$, and so we have
  \[ \Poly_4(\R) = U\oplus W. \qedhere \]
\end{parts}
\end{solution}

\begin{exercise}\leavevmode
\begin{parts}
\item Let $U=\{p\in\Poly_4(\F) : p(2)=p(5)\}$. Find a basis of $U$.
\item Extend the basis in (a) to a basis of $\Poly_4(\F)$.
\item Find a subspace $W$ of $\Poly_4(\F)$ such that $\Poly_4(\F)=U\oplus W$.
\end{parts}
\end{exercise}
\begin{solution}
\solnote{Manual Remark}{Please see remark at the beginning of this section
regarding the basis of $U$.}
\begin{parts}
\item Let $p=a_0+a_1x+a_2x^2+a_3x^3+a_4x^4\in U$. Then we have $p(2)=p(5)$,
  which gives us the equation
  \begin{align*}
  &a_0+2a_1+4a_2+8a_3+16a_4\\
  &\quad= a_0+5a_1+25a_2+125a_3+625a_4.
  \end{align*}
  We can simplify this equation to get
  \[ -3a_1-21a_2-117a_3-609a_4 = 0. \]
  We can solve this equation for $a_1$ in terms of $a_2,a_3,a_4$:
  \[ a_1 = -7a_2-39a_3-203a_4. \]
  Then we can write any polynomial in $U$ as a linear combination of the
  following four polynomials:
  \begin{gather*}
  u_1=1, \quad u_2=-7x+x^2,\\
  u_3=-39x+x^3, \quad u_4=-203x+x^4.
  \end{gather*}
\item We can extend the basis $\{u_1,u_2,u_3,u_4\}$ of $U$ to a basis of
  $\Poly_4(\F)$ by adding one more polynomial that is linearly independent of
  $u_1,u_2,u_3,u_4$. For example, we can add the polynomial
  \[ u_5 = x, \]
  which is linearly independent of $u_1,u_2,u_3,u_4$: since only
  $u_2,u_3,u_4$ have $x$ terms but contain other higher-degree terms, $u_5$
  cannot be expressed as a linear combination of $u_1,u_2,u_3,u_4$. Thus,
  $\{u_1,u_2,u_3,u_4,u_5\}$ is a basis of $\Poly_4(\F)$.
\item Let $W=\spn u_5=\spn x$. Then we have
  \[ \Poly_4(\F) = U+W, \]
  since every polynomial in $\Poly_4(\F)$ can be written as a sum of a
  polynomial in $U$ and a polynomial in $W$. Moreover, if $p\in U\cap W$, then
  $p=\alpha x$ for some $\alpha\in\F$. Since $p(2)=p(5)$, we have
  $2\alpha=5\alpha$, which implies that $\alpha=0$. Thus, $U\cap W=\{0\}$, and
  so we have
  \[ \Poly_4(\F) = U\oplus W. \qedhere \]
\end{parts}
\end{solution}

\begin{exercise}\leavevmode
\begin{parts}
\item Let $U=\{p\in\Poly_4(\F) : p(2)=p(5)=p(6)\}$. Find a basis of $U$.
\item Extend the basis in (a) to a basis of $\Poly_4(\F)$.
\item Find a subspace $W$ of $\Poly_4(\F)$ such that $\Poly_4(\F)=U\oplus W$.
\end{parts}
\end{exercise}
\begin{solution}
\solnote{Manual Remark}{Please see remark at the beginning of this section
regarding the basis of $U$.}
\begin{parts}
\item Let $p=a_0+a_1x+a_2x^2+a_3x^3+a_4x^4\in U$. Then we have
  $p(2)=p(5)=p(6)$. We know from Exercise~2C5 that the condition $p(2)=p(5)$
  gives us the equation
  \[ a_1 = -7a_2-39a_3-203a_4. \]
  The condition $p(5)=p(6)$ gives us the equation
  \begin{align*}
  &5a_1+25a_2+125a_3+625a_4\\
  &\quad= 6a_1+36a_2+216a_3+1296a_4,
  \end{align*}
  which simplifies to
  \[ a_1+11a_2+91a_3+671a_4 = 0. \]
  Since we know that $a_1=-7a_2-39a_3-203a_4$, we can substitute this into the
  second equation to get
  \[ 4a_2+52a_3+468a_4 = 0. \]
  We can solve this equation for $a_2$ in terms of $a_3,a_4$:
  \[ a_2 = -13a_3-117a_4. \]
  Then $a_1=52a_3+616a_4$. Thus, we can write any polynomial in $U$ as a linear
  combination of the following three polynomials:
  \begin{gather*}
  u_1=1, \quad u_2=52x-13x^2+x^3,\\
  u_3=616x-117x^2+x^4.
  \end{gather*}
\item We can extend the basis $\{u_1,u_2,u_3\}$ of $U$ to a basis of
  $\Poly_4(\F)$ by adding two more polynomials that are linearly independent of
  $u_1,u_2,u_3$. For example, we can add the polynomials
  \[ u_4=x^2, \quad u_5=x^3, \]
  which are linearly independent of $u_1,u_2,u_3$: since only $u_2$ and $u_3$
  have $x^2$ terms but contain other higher-degree terms, $u_4$ cannot be
  expressed as a linear combination of $u_1,u_2,u_3$. Similarly, $u_5$ has an
  $x^3$ term but no higher-degree terms, so it cannot be expressed as a linear
  combination of $u_1,u_2,u_3,u_4$. Thus, $\{u_1,u_2,u_3,u_4,u_5\}$ is a basis
  of $\Poly_4(\F)$.
\item Let $W=\spn(u_4,u_5)=\spn(x^2,x^3)$. Then we have
  \[ \Poly_4(\F) = U+W, \]
  since every polynomial in $\Poly_4(\F)$ can be written as a sum of a
  polynomial in $U$ and a polynomial in $W$. Moreover, if $p\in U\cap W$, then
  $p=\alpha x^2+\beta x^3$ for some $\alpha,\beta\in\F$. Since
  $p(2)=p(5)=p(6)$, we have
  \[ 4\alpha+8\beta = 25\alpha+125\beta = 36\alpha+216\beta. \]
  Solving this system of equations, we find that $\alpha=\beta=0$. Thus,
  $U\cap W=\{0\}$, and so we have
  \[ \Poly_4(\F) = U\oplus W. \qedhere \]
\end{parts}
\end{solution}

\begin{exercise}\leavevmode
\begin{parts}
\item Let $U=\bigl\{p\in\Poly_4(\R) : \int_{-1}^1p=0\bigr\}$. Find a basis
  of $U$.
\item Extend the basis in (a) to a basis of $\Poly_4(\R)$.
\item Find a subspace $W$ of $\Poly_4(\R)$ such that $\Poly_4(\R)=U\oplus W$.
\end{parts}
\end{exercise}
\begin{solution}
\solnote{Manual Remark}{Please see remark at the beginning of this section
regarding the basis of $U$.}
\begin{parts}
\item Let $p=a_0+a_1x+a_2x^2+a_3x^3+a_4x^4\in U$. Then we have
  \begin{align*}
  \int_{-1}^1p &= \int_{-1}^1(a_0+a_1x+a_2x^2+a_3x^3+a_4x^4)\,\mathrm{d}x\\
    &= \int_{-1}^1(a_0+a_2x^2+a_4x^4)\,\mathrm{d}x\\
    &= 2\int_0^1(a_0+a_2x^2+a_4x^4)\,\mathrm{d}x\\
    &= 2\Bigl(a_0+\frac{a_2}{3}+\frac{a_4}{5}\Bigr) = 0.
  \end{align*}
  We can solve this equation for $a_0$ in terms of $a_2,a_4$:
  \[ a_0 = -\frac{a_2}{3}-\frac{a_4}{5}. \]
  Then we can write any polynomial in $U$ as a linear combination of the
  following four polynomials:
  \[ u_1=x, \quad u_2=x^2-1/3, \quad u_3=x^3, \quad u_4=x^4-1/5. \]
\item We can extend the basis $\{u_1,u_2,u_3,u_4\}$ of $U$ to a basis of
  $\Poly_4(\R)$ by adding one more polynomial that is linearly independent of
  $u_1,u_2,u_3,u_4$. For example, we can add the polynomial
  \[ u_5 = 1, \]
  which is linearly independent of $u_1,u_2,u_3,u_4$: since $u_5$ has no
  higher-degree terms, it cannot be expressed as a linear combination of
  $u_1,u_2,u_3,u_4$. Thus, $\{u_1,u_2,u_3,u_4,u_5\}$ is a basis of
  $\Poly_4(\R)$.
\item Let $W=\spn u_5=\spn 1$. Then we have
  \[ \Poly_4(\R) = U+W, \]
  since every polynomial in $\Poly_4(\R)$ can be written as a sum of a
  polynomial in $U$ and a polynomial in $W$. Moreover, if $p\in U\cap W$, then
  $p=\alpha$ for some $\alpha\in\R$. Since the integral of a constant over a
  symmetric interval is nonzero unless the function is identically zero, we
  have $\alpha=0$. Thus, $U\cap W=\{0\}$, and so we have
  \[ \Poly_4(\R) = U\oplus W. \qedhere \]
\end{parts}
\end{solution}

\begin{exercise}
Suppose $v_1,\dots,v_m$ is linearly independent in $V$ and $w\in V$. Prove that
\[ \dim\spn(v_1+w,\dots,v_m+w) \ge m-1. \]
\end{exercise}
\begin{solution}
Observe that $(v_i+w)-(v_{i+1}+w)=v_i-v_{i+1}$ for $1\le i\le m-1$. Then the
list of vectors
\[ v_1-v_2,\,v_2-v_3,\,\dots,\,v_{m-1}-v_m \]
is contained in $\spn(v_1+w,\dots,v_m+w)$. We claim that this list is linearly
independent. Suppose that
\[ \alpha_1(v_1-v_2)+\alpha_2(v_2-v_3)+\cdots+\alpha_{m-1}(v_{m-1}-v_m) = 0 \]
for some $\alpha_1,\dots,\alpha_{m-1}\in\F$. Then we can rewrite this equation
as
\begin{align*}
&\alpha_1v_1+(\alpha_2-\alpha_1)v_2+\cdots\\
&\quad+(\alpha_{m-1}-\alpha_{m-2})v_{m-1}-\alpha_{m-1}v_m = 0.
\end{align*}
Since $v_1,\dots,v_m$ is linearly independent, we must have $\alpha_1=0$,
$\alpha_2-\alpha_1=0$, \dots, $\alpha_{m-1}-\alpha_{m-2}=0$, and
$-\alpha_{m-1}=0$. This implies that $\alpha_1=\cdots=\alpha_{m-1}=0$. Thus, the
list is linearly independent, and so we have
\[ \dim\spn(v_1+w,\dots,v_m+w) \ge m-1. \]

\emph{Alternate proof.} \emph{Civil Service Pigeon.} Let
$S=\spn(v_1+w,\dots,v_m+w)$. Since each $v_i=(v_i+w)-w$, we have
\[ S+\spn(w) = \spn(v_1,\dots,v_m,w). \]
Since $v_1,\dots,v_m$ are linearly independent, we have
\begin{align*}
m &\le \dim\bigl(S+\spn(w)\bigr) \le \dim S+\dim\spn(w)\\
  &\le \dim S+1 \implies \dim S\ge m-1. \qedhere
\end{align*}
\end{solution}

\begin{exercise}
Suppose $m$ is a positive integer and $p_0,p_1,\dots,p_m\in\Poly(\F)$ are such
that each $p_k$ has degree $k$. Prove that $p_0,p_1,\dots,p_m$ is a basis of
$\Poly_m(\F)$.
\end{exercise}
\begin{solution}
Suppose that
\[ \alpha_0p_0+\alpha_1p_1+\cdots+\alpha_mp_m = 0 \]
for some $\alpha_0,\dots,\alpha_m\in\F$. Since $p_m$ is the only polynomial with
nonzero coefficient of $x^m$, we must have $\alpha_m=0$. Next, since $p_{m-1}$
is the only polynomial with nonzero coefficient of $x^{m-1}$, we must have
$\alpha_{m-1}=0$. Continuing in this way, we find that
$\alpha_0=\cdots=\alpha_m=0$. Thus, the list is linearly independent. Moreover,
since $\dim\Poly_m(\F)=m+1$, any linearly independent list of $m+1$ vectors in
$\Poly_m(\F)$ is a basis. Therefore, $p_0,p_1,\dots,p_m$ is a basis of
$\Poly_m(\F)$.
\end{solution}

\begin{exercise}
Suppose $m$ is a positive integer. For $0\le k\le m$, let
\[ p_k(x) = x^k(1-x)^{m-k}. \]
Show that $p_0,\dots,p_m$ is a basis of $\Poly_m(\F)$.
\exnote{The basis in this exercise leads to what are called \textbf{Bernstein
polynomials}. You can do a web search to learn how Bernstein polynomials are
used to approximate continuous functions on $[0,1]$.}
\end{exercise}
\begin{solution}
Suppose that
\[ \alpha_0p_0+\alpha_1p_1+\cdots+\alpha_mp_m = 0 \]
for some $\alpha_0,\dots,\alpha_m\in\F$. Since $p_0$ is the only polynomial with
nonzero constant term, we must have $\alpha_0=0$, hence the sum above becomes
\[ \alpha_1p_1+\cdots+\alpha_mp_m = 0. \]
Next, since $p_1$ is the only polynomial with nonzero coefficient of $x$ in the
sum above, we must have $\alpha_1=0$. Continuing in this way, we find that
$\alpha_0=\cdots=\alpha_m=0$. Thus, the list is linearly independent. Moreover,
since $\dim\Poly_m(\F)=m+1$, any linearly independent list of $m+1$ vectors in
$\Poly_m(\F)$ is a basis. Therefore, $p_0,\dots,p_m$ is a basis of
$\Poly_m(\F)$.
\end{solution}

\begin{exercise}
Suppose $U$ and $W$ are both four-dimensional subspaces of $\C^6$. Prove that
there exist two vectors in $U\cap W$ such that neither of these vectors is a
scalar multiple of the other.
\end{exercise}
\begin{solution}
Let $U$ and $W$ be four-dimensional subspaces of $\C^6$. By the dimension
formula for the sum of two subspaces, we have
\begin{align*}
\dim(U+W) &= \dim U+\dim W-\dim(U\cap W)\\
  &= 4+4-\dim(U\cap W) = 8-\dim(U\cap W).
\end{align*}
Since $U+W$ is a subspace of $\C^6$, we have $\dim(U+W)\le6$. Thus, we have
\[ 8-\dim(U\cap W) \le 6, \]
which implies that $\dim(U\cap W)\ge2$. Therefore, there exist at least two
linearly independent vectors in $U\cap W$, and neither of these vectors is a
scalar multiple of the other.
\end{solution}

\begin{exercise}
Suppose that $U$ and $W$ are subspaces of $\R^8$ such that $\dim U=3$,
$\dim W=5$, and $U+W=\R^8$. Prove that $\R^8=U\oplus W$.
\end{exercise}
\begin{solution}
By the dimension formula for the sum of two subspaces, we have
\begin{align*}
\dim(U+W) &= \dim U+\dim W-\dim(U\cap W)\\
  &= 3+5-\dim(U\cap W) = 8-\dim(U\cap W).
\end{align*}
Since $U+W=\R^8$, we have $\dim(U+W)=8$. Thus, we have
\[ 8-\dim(U\cap W) = 8, \]
which implies that $\dim(U\cap W)=0$. Therefore, $U\cap W=\{0\}$, and so we
have
\[ \R^8 = U\oplus W. \qedhere \]
\end{solution}

\begin{exercise}
Suppose $U$ and $W$ are both five-dimensional subspaces of $\R^9$. Prove that
$U\cap W\ne\{0\}$.
\end{exercise}
\begin{solution}
By the dimension formula for the sum of two subspaces, we have
\begin{align*}
\dim(U+W) &= \dim U+\dim W-\dim(U\cap W)\\
  &= 5+5-\dim(U\cap W) = 10-\dim(U\cap W).
\end{align*}
Since $U+W$ is a subspace of $\R^8$, we have $\dim(U+W)\le8$. Thus, we have
\[ 10-\dim(U\cap W) \le 8, \]
which implies that $\dim(U\cap W)\ge2$. Therefore, $U\cap W$ is nontrivial, and
so we have $U\cap W\ne\{0\}$.
\end{solution}

\begin{exercise}
Suppose $V$ is a ten-dimensional vector space and $V_1,V_2,V_3$ are subspaces
of $V$ with $\dim V_1=\dim V_2=\dim V_3=7$. Prove that
$V_1\cap V_2\cap V_3\ne\{0\}$.
\end{exercise}
\begin{solution}
Let $W=V_1\cap V_2$. By the dimension formula for the sum of two subspaces, we
have
\begin{align*}
\dim(V_1+V_2) &= \dim V_1+\dim V_2-\dim W\\
  &= 7+7-\dim W = 14-\dim W.
\end{align*}
Since $V_1+V_2$ is a subspace of $V$, we have $\dim(V_1+V_2)\le10$. Thus, we
have $14-\dim W\le10$, which implies that $\dim W\ge4$. Therefore, $V_1\cap V_2$
is at least four-dimensional. Now, consider the subspace
$W\cap V_3=V_1\cap V_2\cap V_3$. By the dimension formula for the sum of two
subspaces, we have
\begin{align*}
\dim(W+V_3) &= \dim W+\dim V_3-\dim(W\cap V_3)\\
  &= \dim W+7-\dim(W\cap V_3).
\end{align*}
Since $W+V_3$ is a subspace of $V$, we have $\dim(W+V_3)\le10$. Thus, we have
\[ \dim W+7-\dim(W\cap V_3) \le 10, \]
which implies that $\dim(W\cap V_3)\ge\dim W+7-10=\dim W-3$. Since
$\dim W\ge4$, we have $\dim(W\cap V_3)\ge1$. Therefore, $V_1\cap V_2\cap V_3$ is
nontrivial, and so we have $V_1\cap V_2\cap V_3\ne\{0\}$.
\end{solution}

\begin{exercise}
Suppose $V$ is finite-dimensional and $V_1,V_2,V_3$ are subspaces of $V$ with
$\dim V_1+\dim V_2+\dim V_3>2\dim V$. Prove that $V_1\cap V_2\cap V_3\ne\{0\}$.
\end{exercise}
\begin{solution}
Let $W=V_1\cap V_2$. By the dimension formula for the sum of two subspaces, we
have
\[ \dim(V_1+V_2) = \dim V_1+\dim V_2-\dim W. \]
Since $V_1+V_2$ is a subspace of $V$, we have $\dim(V_1+V_2)\le\dim V$. Thus, we
have
\[ \dim V_1+\dim V_2-\dim W \le \dim V, \]
which implies that $\dim W\ge\dim V_1+\dim V_2-\dim V$. Now, consider the
subspace $W\cap V_3=V_1\cap V_2\cap V_3$. By the dimension formula for the sum
of two subspaces, we have
\[ \dim(W+V_3) = \dim W+\dim V_3-\dim(W\cap V_3). \]
Since $W+V_3$ is a subspace of $V$, we have $\dim(W+V_3)\le\dim V$. Thus, we
have
\[ \dim W+\dim V_3-\dim(W\cap V_3) \le \dim V, \]
which implies that $\dim(W\cap V_3)\ge\dim W+\dim V_3-\dim V$. Since
$\dim W\ge\dim V_1+\dim V_2-\dim V$, we have
\begin{align*}
\dim(W\cap V_3) &\ge (\dim V_1+\dim V_2-\dim V)+\dim V_3-\dim V\\
  &= \dim V_1+\dim V_2+\dim V_3-2\dim V.
\end{align*}
Since $\dim V_1+\dim V_2+\dim V_3>2\dim V$, we have $\dim(W\cap V_3)>0$.
Therefore, $V_1\cap V_2\cap V_3$ is nontrivial, and so we have
$V_1\cap V_2\cap V_3\ne\{0\}$.
\end{solution}

\begin{exercise}
Suppose $V$ is finite-dimensional and $U$ is a subspace of $V$ with $U\ne V$.
Let $n=\dim V$ and $m=\dim U$. Prove that there exist $n-m$ subspaces of $V$,
each of dimension $n-1$, whose intersection equals $U$.
\end{exercise}
\begin{solution}
Let $U$ be a subspace of $V$ with $\dim U=m<n=\dim V$. We can extend a basis of
$U$ to a basis of $V$. Let $\{u_1,\dots,u_m\}$ be a basis of $U$, and extend it
to a basis of $V$ by adding vectors $v_1,\dots,v_{n-m}$. Then
$\{u_1,\dots,u_m,v_1,\dots,v_{n-m}\}$ is a basis of $V$. For each
$i=1,\dots,n-m$, let
\[ W_i = \spn(u_1,\dots,u_m,v_1,\dots,v_{i-1},v_{i+1},\dots,v_{n-m}). \]
By construction, $W_i$ does not contain the extended basis vector $v_i$.
Moreover, each $W_i$ is a subspace of $V$ with dimension $n-1$, since it is
spanned by $n-1$ linearly independent vectors. Additionally, we have
\[ U = W_1\cap W_2\cap\cdots\cap W_{n-m}, \]
since any vector in the intersection containing nonzero $v_i$ would not be in
$W_i$, leaving behind only the basis vectors in $U$. Therefore, there exist
$n-m$ subspaces of $V$, each of dimension $n-1$, whose intersection equals $U$.
\end{solution}

\begin{exercise}
Suppose that $V_1,\dots,V_m$ are finite-dimensional subspaces of $V$. Prove that
$V_1+\cdots+V_m$ is finite-dimensional and
\[ \dim(V_1+\cdots+V_m) \le \dim V_1+\cdots+\dim V_m. \]
\exnote{The inequality above is an equality if and only if $V_1+\cdots+V_m$ is
a direct sum, as will be shown in \ref{ax:3.94}.}
\end{exercise}
\begin{solution}
We will prove the statement by induction on $m$. For the base case, let $m=2$.
Then the sum $V_1+V_2$ is finite-dimensional. By the dimension formula for the
sum of two subspaces, we have
\begin{align*}
\dim(V_1+V_2) &= \dim V_1+\dim V_2-\dim(V_1\cap V_2)\\
  &\le \dim V_1+\dim V_2.
\end{align*}
Hence, the statement holds for $m=2$. Now, suppose that the statement holds for
some positive integer $m$. That is, suppose that if $V_1,\dots,V_m$ are
finite-dimensional subspaces of $V$, then $V_1+\cdots+V_m$ is finite-dimensional
and
\[ \dim(V_1+\cdots+V_m) \le \dim V_1+\cdots+\dim V_m. \]
We want to show that the statement holds for $m+1$. Let $V_1,\dots,V_{m+1}$ be
finite-dimensional subspaces of $V$. By the induction hypothesis, we know that
$W=V_1+\cdots+V_m$ is finite-dimensional and
\[ \dim W \le \dim V_1+\cdots+\dim V_m. \]
Then we have
\[ \dim(W+V_{m+1}) = \dim W+\dim V_{m+1}-\dim(W\cap V_{m+1}). \]
Since $W\cap V_{m+1}$ is a subspace of $V$, we have $\dim(W\cap V_{m+1})\ge0$.
Thus, we have
\begin{align*}
\dim(W+V_{m+1}) &\le \dim W+\dim V_{m+1}\\
  &\le \dim V_1+\cdots+\dim V_m+\dim V_{m+1}.
\end{align*}
Therefore, $V_1+\cdots+V_{m+1}$ is finite-dimensional and
\[ \dim(V_1+\cdots+V_{m+1}) \le \dim V_1+\cdots+\dim V_m+\dim V_{m+1}.
   \qedhere \]
\end{solution}

\begin{exercise}
Suppose $V$ is finite-dimensional, with $\dim V=n\ge1$. Prove that there exist
one-dimensional subspaces $V_1,\dots,V_n$ of $V$ such that
\[ V = V_1\oplus\cdots\oplus V_n. \]
\end{exercise}
\begin{solution}
Let $v_1,\dots,v_n$ be a basis of $V$. For each $i=1,\dots,n$, let
$V_i=\spn v_i$. Then each $V_i$ is a one-dimensional subspace of $V$, since it
is spanned by a single nonzero vector. Moreover, we have
\[ V = V_1+\cdots+V_n, \]
and the sum is direct because the basis vectors are linearly independent.
Therefore, we have
\[ V = V_1\oplus\cdots\oplus V_n. \qedhere \]
\end{solution}

\begin{exercise}
Explain why you might guess, motivated by analogy with the formula for the
number of elements in the union of three finite sets, that if $V_1,V_2,V_3$ are
subspaces of a finite-dimensional vector space, then
\begin{align*}
\dim(V_1&+V_2+V_3)\\
  &= \dim V_1+\dim V_2+\dim V_3\\
  &\quad-\dim(V_1\cap V_2)-\dim(V_1\cap V_3)-\dim(V_2\cap V_3)\\
  &\quad+\dim(V_1\cap V_2\cap V_3).
\end{align*}
Then either prove the formula above or give a counterexample.
\end{exercise}
\begin{solution}
One may guess that the given formula above is true by analogy with the formula
for the number of elements in the union of three finite sets. However, this
erroneously conflates the idea of subspace sums with the idea of set unions. In
fact, the formula is not true in general. For example, let $V=\R^2$, and let
\[ V_1=\spn\bigl((1,0)\bigr), \quad V_2=\spn\bigl((0,1)\bigr), \quad
   V_3=\spn\bigl((1,1)\bigr). \]
It is easy to see that $\dim V_1=\dim V_2=\dim V_3=1$, but
$\dim(V_1\cap V_2)=\dim(V_1\cap V_3)=\dim(V_2\cap V_3)=0$ and
$\dim(V_1\cap V_2\cap V_3)=0$. Thus, the right-hand side of the formula is
\[ 1+1+1-0-0-0+0 = 3, \]
but the left-hand side is $\dim(V_1+V_2+V_3)=\dim\R^2=2$. Therefore, the
formula is not true in general.
\end{solution}

\begin{exercise}
Prove that if $V_1$, $V_2$, and $V_3$ are subspaces of a finite-dimensional
vector space, then
\begin{align*}
&\dim(V_1+V_2+V_3)\\
&\quad= \dim V_1+\dim V_2+\dim V_3\\
&\qquad-\frac{\dim(V_1\cap V_2)+\dim(V_1\cap V_3)+\dim(V_2\cap V_3)}{3}\\
&\qquad-\frac{\dim\bigl((V_1+V_2)\cap V_3\bigr)+\dim\bigl((V_1+V_3)\cap V_2\bigr)
  +\dim\bigl((V_2+V_3)\cap V_1\bigr)}{3}.
\end{align*}
\exnote{The formula above may seem strange because the right side does not look
like an integer.}
\end{exercise}
\begin{solution}
Note that we need to account for all pairs of subspaces and all pairs of sums of
subspaces, since the sum of subspaces is not equivalent to the union of sets. We
can use the dimension formula repeatedly to obtain the answer. To that end,
define the following subspaces:
\[ W_1=V_1+V_2, \quad W_2=V_1+V_3, \quad W_3=V_2+V_3. \]
Applying the dimension formula to each of $W_1,W_2,W_3$, we have
\begin{align*}
\dim W_1 &= \dim V_1+\dim V_2-\dim(V_1\cap V_2),\\
\dim W_2 &= \dim V_1+\dim V_3-\dim(V_1\cap V_3),\\
\dim W_3 &= \dim V_2+\dim V_3-\dim(V_2\cap V_3).
\end{align*}
Next, we can apply the dimension formula to each of the sums $W_1+V_3$,
$W_2+V_2$, $W_3+V_1$:
\begin{align*}
\dim(W_1+V_3) &= \dim W_1+\dim V_3-\dim(W_1\cap V_3),\\
  &= \dim V_1+\dim V_2-\dim(V_1\cap V_2)\\
  &\qquad+\dim V_3-\dim(W_1\cap V_3)\\
\dim(W_2+V_2) &= \dim W_2+\dim V_2-\dim(W_2\cap V_2),\\
  &= \dim V_1+\dim V_3-\dim(V_1\cap V_3)\\
  &\qquad+\dim V_2-\dim(W_2\cap V_2)\\
\dim(W_3+V_1) &= \dim W_3+\dim V_1-\dim(W_3\cap V_1),\\
  &= \dim V_2+\dim V_3-\dim(V_2\cap V_3)\\
  &\qquad+\dim V_1-\dim(W_3\cap V_1).
\end{align*}
Note that $W_1+V_3=W_2+V_2=W_3+V_1=V_1+V_2+V_3$. Thus, we can add the three
equations above and divide by $3$ to obtain the desired formula:
\begin{align*}
&\dim(V_1+V_2+V_3)\\
&\quad= \dim V_1+\dim V_2+\dim V_3\\
&\qquad-\frac{\dim(V_1\cap V_2)+\dim(V_1\cap V_3)+\dim(V_2\cap V_3)}{3}\\
&\qquad-\frac{\dim\bigl((V_1+V_2)\cap V_3\bigr)+\dim\bigl((V_1+V_3)\cap V_2\bigr)
  +\dim\bigl((V_2+V_3)\cap V_1\bigr)}{3}. \qedhere
\end{align*}
\end{solution}
\end{exc}

\axquote{I at once gave up my former occupations, set down natural history and
all its progeny as a deformed and abortive creation, and entertained the
greatest disdain for a would-be science which could never even step within the
threshold of real knowledge. In this mood I betook myself to the mathematics
and the branches of study appertaining to that science as being built upon
secure foundations, and so worthy of my consideration.}{\emph{Frankenstein},
Mary Wollstonecraft Shelley}
